\clearpage
\SourceLeaf{169}%
\pdfbookmark[0]{Göttingische gelehrte Anzeigen. 1817 March 10.}{notice1817march10}
\begin{center}
\rule{.82\linewidth}{.35pt}\par
Göttingische gelehrte Anzeigen.\quad 1817 March 10.\par
\rule{.82\linewidth}{.35pt}
\end{center}
On 10 February Court Councillor \textsc{Gauss} submitted a lecture to the Royal Society, entitled:
\begin{center}\itshape
Theorematis fundamentalis in doctrina de residuis quadraticis demonstrationes et ampliationes novae.
\end{center}

A distinctive feature of higher arithmetic is that so many of its most beautiful theorems can be discovered with the greatest ease by examining examples, while their proofs are far from readily apparent. Often they are found only after many unsuccessful attempts, through deeply probing investigations and fortunate combinations. This remarkable phenomenon arises from the often wonderful interconnection between the different branches of this part of mathematics. For the same reason, theorems for which a proof was initially sought in vain for years can often later be proved in several entirely different ways. Once a new theorem has been discovered through examples, finding \emph{some} proof must of course be the first requirement; but after this has succeeded, in higher arithmetic we must not always regard the investigation as finished and the search for other proofs as a needless luxury. For one thing, we do not usually find the most beautiful and simplest
\SourceLeaf{170}%
proofs first; for another, insight into the wonderful interconnection of the truth of higher arithmetic is precisely one of the chief attractions of this study, and not infrequently leads in turn to the discovery of new truths. For these reasons, finding new proofs of truths already known must here often be considered at least as important as discovering the truth itself. These reflections are familiar to those who know higher arithmetic: much of \textsc{Euler}'s contribution to it is known to consist in finding proofs of theorems which \textsc{Fermat} had already discovered, apparently through examples.

The theory of quadratic residues provides a clear illustration of these remarks. It rests chiefly on the so-called fundamental theorem, which states that the mutual relations of two (odd positive) primes to one another (insofar as one is a quadratic residue or nonresidue of the other) are the same whenever one or both primes have the form $4k+1$, but opposite whenever both primes have the form $4k+3$. For readers less familiar with higher arithmetic, we recall that an integer is called a quadratic residue of another if adding a multiple of the latter to the former can give a square; it is called a nonresidue if this is impossible. We do not intend to repeat the whole history of this beautiful theorem, so easily discovered through examples. We merely observe that the author of the present paper, after investigations initially pursued in vain for quite some time, has already given, one after another, four entirely different proofs. Two are contained in the \emph{Disquisitionibus Arithmeticis}; the third is the subject of a separate paper in the sixteenth volume of the Commentationes; and the fourth is interwoven with the paper \emph{summatio quarumdam serierum singularium} in the first volume of the \emph{Commentationes recentiores}. For these two papers readers may consult our reviews of 1808, May 12 and September 19, which also give fuller historical references. That the author did not stop at these four proofs, although each leaves nothing to be desired as regards rigour, needs no justification among the friends of higher arithmetic. Still, he would probably not have worked so eagerly to add further proofs if
\SourceLeaf{171}%
a particular circumstance had not prompted him to do so, and this must be mentioned here. In the year 1805 he began to study the theories of cubic and biquadratic residues, which are much richer and more interesting still than the theory of quadratic residues. These investigations revealed the same phenomena as the latter theory, only, as it were, on a larger scale. As soon as the right path was taken, examining examples immediately yielded a number of extremely simple theorems which exhaust those theories, bear a surprising resemblance to the theorems for quadratic residues, and in particular also provide a counterpart to the fundamental theorem. But the difficulty of finding entirely satisfactory proofs of these theorems proved greater still, and only after many attempts, continued over a considerable number of years, did the author finally reach his goal. The close analogy between the theorems for quadratic and higher residues suggested that there must also be analogous proofs for the two. Yet the methods of proof first found for quadratic residues could not be applied to the higher residues at all; this was precisely what led him to keep seeking further new proofs for the former. The author therefore wishes the present paper, which opens up some further resources for the theory of quadratic residues, to be regarded as a precursor to the theory of cubic and biquadratic residues. He intends to publish that theory in the future; its subjects are among the most difficult in higher arithmetic.

The present paper consists of three independent parts. It contains the fifth and sixth proofs of the fundamental theorem, and a new method, connected with the third proof,\textperiodcentered\ for deciding whether a given integer is a quadratic residue or nonresidue of a given prime. Of the first four proofs, the third was unquestionably the one that combined the greatest simplicity with independence from investigations of a different kind; \textsc{Legendre} therefore also included it in the new edition of his \emph{Essai d'une théorie des nombres}. The \emph{fifth} proof seems at least to equal the third in both respects. The two proofs are related insofar as they start from the same theorem, but in their subsequent development they are completely different from \SourceLeaf{172}%
one another. \Needspace{12\baselineskip}The theorem is as follows. Let $m$ be a (positive odd) prime and $M$ an integer not divisible by $m$. If, among the remainders obtained by dividing the products
\[
M,\ 2M,\ 3M,\ 4M\mathinner{.\,.\,.\,.\,.\,.}\tfrac12(m-1)M
\]
by $m$, the number of those greater than $\frac12m$ is denoted by $n$, then $M$ is a quadratic residue or nonresidue of $m$ according as $n$ is even or odd. To obtain a proof of the fundamental theorem, it is now assumed that $M$, too, is an odd positive prime, and that $N$ has the same meaning with respect to $M$ and $m$ as $n$ has with respect to $m$ and $M$. Thus whether $N$ is even or odd determines whether $m$ is a quadratic residue or nonresidue of $M$. By a very short chain of arguments the author shows that the number of all positive integers which are less than $\frac12mM$ and at the same time on division by $m$ leave a remainder less than $\frac12m$, and on division by $M$ a remainder less than $\frac12M$,
\[
=\tfrac18(m-1)(M-1)+\tfrac12n+\tfrac12N
\]
and consequently that
\[
\tfrac14(m-1)(M-1)+n+N
\]
is always an even number. Thus, whenever at least one of the numbers $m,M$ has the form $4k+1$, so that $\frac14(m-1)(M-1)$ is even, $n+N$ will also be even; hence $n$ and $N$ are either both even or both odd. If, on the other hand, both $m$ and $M$ have the form $4k+3$, then $n+N$ must be odd, so one of the numbers $n,N$ is even and the other odd. Together with the theorem above, this gives the fundamental theorem at once.

The \emph{sixth} proof is as short and elegant as the fifth, but rests on somewhat more intricate combinations. The limited space of these pages allows us only to touch on its main point, omitting the details. Let the following symbols denote:
\begin{description}
\item[$p,q$] two (distinct positive odd) primes,
\item[$\alpha$] a so-called \emph{primitive root} for the modulus $p$, that is, an integer not divisible by $p$ (here taken positive), such that no power lower than $\alpha^{p-1}$ is congruent to one modulo $p$
\item[$x$] an indeterminate quantity
\SourceLeaf{173}%
\item[$\xi$] the function
\[
x-x^\alpha+x^\zeta-x^\eta+x^\theta-\text{etc.}-x^\lambda
\]
where, for convenience in printing, $\zeta,\eta,\theta\ldots\lambda$ stand for the numbers $\alpha\alpha,\allowbreak \alpha^3,\allowbreak \alpha^4\ldots\allowbreak \alpha^{p-2}$;
\item[$\varepsilon$] one, taken positively if $p$ has the form $4k+1$, negatively if $p$ has the form $4k+3$;
\item[$\delta$] one, taken positively if at least one of the numbers $p,q$ has the form $4k+1$, negatively if both have the form $4k+3$;
\item[$\gamma$] one, taken positively if $q$ is a quadratic residue of $p$, negatively if $q$ is a quadratic nonresidue of $p$;
\item[$\beta$] one, taken positively if $p$ is a quadratic residue of $q$, negatively if $p$ is a quadratic nonresidue of $q$.
\end{description}

With these preparations it follows readily from Article 51 of the \emph{Disquisitiones Arithmeticae} that the function
\[
\xi^q-x^q+x^{q\alpha}-x^{q\zeta}+x^{q\eta}-x^{q\theta}+\text{etc.}+x^{q\lambda}
\]
has, when expanded, coefficients all divisible by $q$. Hence, if this function is set $=qX$, then $X$ is a \emph{polynomial} whose coefficients, too, are integers. By arguments which it would take too much space to discuss here, the paper proves that dividing the function $qX\xi$ by $x^{p-1}+x^{p-2}+x^{p-3}+x^{p-4}+\text{etc.}+x+1$ leaves the remainder
\[
\varepsilon p(\delta p^{\frac12(q-1)}-\gamma)
\]
and therefore that dividing the function $X\xi$ by the same divisor leaves the remainder
\[
\frac{\varepsilon p(\delta p^{\frac12(q-1)}-\gamma)}{q}
\]
This quantity must therefore be an integer. Since $\delta\delta=1$, it follows readily that
\[
p^{\frac12(q-1)}-\gamma\delta
\]
must be divisible by $q$. A known theorem also gives that $p^{\frac12(q-1)}-\beta$ is divisible by $q$, so necessarily $\beta=\gamma\delta$, from which the fundamental theorem again follows at once.

\clearpage\SourceLeaf{174}%
The fundamental theorem, together with some known theorems, can indeed provide a fairly short way of deciding whether a given positive integer is a quadratic residue or nonresidue of a given prime, as the paper shows in detail. But further reflection on the third proof of the fundamental theorem led the author to a much more flexible solution. This forms the paper's third part; here we give only its final rule, omitting the reasoning behind it for brevity. To decide whether the positive integer $b$, not divisible by the prime $a$, is a quadratic residue or nonresidue of that prime, form, in just the same way as when finding the greatest common divisor of $a$ and $b$, the equations
\[
\begin{aligned}
a&=\beta b+c\\
b&=\gamma c+d\\
c&=\delta d+e\\
d&=\varepsilon e+f\text{ etc.}
\end{aligned}
\]
until the sequence of numbers $a,b,c,d,e,f$ and so on reaches .the unit. Let. the numbers $\frac12a,\frac12b,\frac12c,\frac12d$ and so on, with their fractional part $\frac12$ discarded whenever any of the numbers $a,b,c,d$ and so on is odd, be denoted by $a',b',c',d'$ and so on. Let $\mu$ be the number of occurrences of two odd numbers immediately following one another in the sequence $a',b',c',d'$ and so on. Finally, let $\nu$ be the number of odd numbers in the sequence $\beta,\gamma,\delta,\varepsilon$ and so on to which, in the sequence $b',c',d',e'$ and so on, a number of the form $4k+1$ or $4k+2$ respectively corresponds. Then $b$ is a quadratic residue or nonresidue of $a$ according as $\mu+\nu$ is even or odd, except in the single case where $b$ is even and $a$ has the form $8k+3$ or $8k+5$. In that case the rule is reversed: an even $\mu+\nu$ indicates that $b$ is a quadratic nonresidue of $a$, while an odd $\mu+\nu$ indicates that $b$ is a quadratic residue of $a$.
\BookRule{.28\linewidth}
